

\subsection{Ablation Study}





% \begin{table}[!h]
% \centering
% \setlength{\tabcolsep}{0.1cm}
% \begin{minipage}[t]{.49\linewidth}
% \centering
% \caption{\textbf{Effect of staleness under a fixed generation budget.}
% All settings use the same data generated from a single rollout generation phase.
% The batch size therefore determines both the number of model update steps and
% the resulting staleness level $\mu$. Experiments are conducted with Qwen2.5-Math-7B. }
% \label{tab:batch-mu-ablation}
% \small
% \resizebox{\linewidth}{!}{
% \begin{tabular}{cccccccc}
% \toprule
% \scriptsize Batch Size & \scriptsize $\mu$
% & \scriptsize AMC23 & \scriptsize AMC24
% & \scriptsize AIME24 & \scriptsize AIME25
% & \scriptsize MATH500
% & \scriptsize Avg. \\
% \midrule
% \rowcolor{gray!20}  32  & 1259 & 61.9 & 40.6 & 24.8 & 13.8 & 80.0 & \textbf{44.2} \\
% 64  & 629  & 63.1 & 39.4 & 23.3 & 13.3 & 79.7 & 43.8 \\
% 128 & 314  & 61.3 & 36.1 & 21.5 & 13.5 & 80.7 & 42.6 \\
% 256 & 157  & 62.5 & 32.8 & 23.1 & 11.9 & 79.3 & 41.9 \\
% \bottomrule
% \end{tabular}}    
% \end{minipage}
% \begin{minipage}[t]{.49\linewidth}
% \centering
% \setlength{\tabcolsep}{0.1cm}
% \caption{\textbf{Effect of staleness under a fixed optimization budget.}
% All settings perform 4{,}096 model update steps with a batch size of 32,
% while varying the staleness level $\mu$. Experiments are conducted with Qwen2.5-Math-7B. The default setting has gray background.} 
% \label{tab:staleness_ablation}
% \small
% \resizebox{\linewidth}{!}{
% \begin{tabular}{cccccccc}
% \toprule
% \scriptsize $\mu$
% & \scriptsize AMC23
% & \scriptsize AMC24
% & \scriptsize AIME24
% & \scriptsize AIME25
% & \scriptsize MATH500
% & \scriptsize Avg.
% & \scriptsize \#Stage \\
% \midrule
% 512  & 65.6 & 43.9 & 30.8 & 13.3 & 82.9 & 47.3 & 8 \\
% \rowcolor{gray!20} 1024  & 63.1 & 47.2 & 31.5 & 16.5 & 81.7 & \textbf{48.0} & 4 \\
% 2048 
%   & 65.7 & 40.6 & 28.1 & 16.9 & 79.1 & 46.1 & 2 \\
% \bottomrule
% \end{tabular}}
% \end{minipage}
% \vspace{-10pt}
% \end{table}

 
\begin{table}[!h]
\centering
\setlength{\tabcolsep}{0.08cm}
\small
\begin{minipage}[t]{.48\linewidth}
\centering
\caption{\textbf{Fixed generation budget.}
Varying batch size on a fixed rollout dataset.}
\label{tab:batch-mu-ablation}
\resizebox{\linewidth}{!}{
\begin{tabular}{cccccccc}
\toprule
\scriptsize Batch & \scriptsize $\mu$
& \scriptsize AMC23 & \scriptsize AMC24
& \scriptsize AIME24 & \scriptsize AIME25
& \scriptsize MATH500
& \scriptsize Avg. \\
\midrule
\rowcolor{gray!20} 
32  & 1259 & 61.9 & 40.6 & 24.8 & 13.8 & 80.0 & \textbf{44.2} \\
64  & 629  & 63.1 & 39.4 & 23.3 & 13.3 & 79.7 & 43.8 \\
128 & 314  & 61.3 & 36.1 & 21.5 & 13.5 & 80.7 & 42.6 \\
256 & 157  & 62.5 & 32.8 & 23.1 & 11.9 & 79.3 & 41.9 \\
\bottomrule
\end{tabular}}
\end{minipage}
\hfill
\begin{minipage}[t]{.486\linewidth}
\centering
\caption{\textbf{Fixed optimization budget.}
Varying $\mu$ with 4096 updates.}
\label{tab:staleness_ablation}
\resizebox{\linewidth}{!}{
\begin{tabular}{cccccccc}
\toprule
\scriptsize $\mu$
& \scriptsize AMC23
& \scriptsize AMC24
& \scriptsize AIME24
& \scriptsize AIME25
& \scriptsize MATH500
& \scriptsize Avg.
& \scriptsize \#Stage \\
\midrule
512  & 65.6 & 43.9 & 30.8 & 13.3 & 82.9 & 47.3 & 8 \\
\rowcolor{gray!20}
1024  & 63.1 & 47.2 & 31.5 & 16.5 & 81.7 & \textbf{48.0} & 4 \\
2048 & 65.7 & 40.6 & 28.1 & 16.9 & 79.1 & 46.1 & 2 \\
4096 & 66.9 & 37.2 & 26.0 & 16.0 & 82.3 & 45.7 & 1 \\
\bottomrule
\end{tabular}}
\end{minipage}
\vspace{-10pt}
\end{table}
\paragraph{Effect of Staleness with a Fixed Generation Budget.}
To isolate the value of reusing stale rollout data, we first generate a fixed
rollout dataset by sampling 8 responses for each prompt in the DeepScaleR
training set, and then optimize only on this dataset. Under this setting, the batch size determines both the number of updates and the
effective staleness $\mu$: smaller batches yield more updates but higher
staleness, while larger batches reduce staleness but shorten training.
\cref{tab:batch-mu-ablation} varies the batch size in $\{32,64,128,256\}$.
Batch size 32 ($\mu=1259$) achieves the best average accuracy of 44.2\%, while
batch size 256 drops to 41.9\%. This suggests that rollout data remain useful
for many updates, and that additional updates can outweigh the cost of
higher staleness.


\paragraph{Effect of Staleness with a Fixed Optimization Budget.}
\cref{tab:staleness_ablation} fixes the total number of model updates and varies
the effective staleness $\mu$ by changing the rollout refresh frequency, or
equivalently the number of generation--optimization stages. Performance remains
stable from $\mu\!=\!512$ to $\mu\!=\!2048$, but degrades when using a single stage
($\mu\!=\!4096$). Among these settings, $\mu\!=\!1024$ provides the best average
accuracy and is therefore adopted as the
default configuration.



% Intuitively, low staleness may seem preferable, as the training data remains more representative of the current policy.
% Yet larger batch sizes necessarily imply fewer gradient steps, resulting in coarser optimization that may fail to fully exploit the learning signal contained in the data.
% On the other hand, although smaller batch sizes lead to higher staleness, they enable more frequent and finer-grained policy updates.
% Such fine-grained updates can improve sample efficiency by allowing the optimizer to make incremental adjustments based on smaller subsets of the data, which often leads to better generalization and more stable convergence.



